% !TeX root = ../../Book5.tex
% 纲要标题：### 19.5 维数与重数
% 纲要位置：工作记录/计划纲要.md，第 4342 行。
% * Weyl 维数公式
% * 权重数的计算
% * 特征标乘法与张量分解
% * 由形式代换或最低次项比较导出 Weyl 维数公式；不把分子分母同时为零时的直接代入当作证明

\section{维数与重数}
\label{b5:sec:1905}

Weyl 特征标公式已经同时包含维数和全部权重数，但取出这两类信息的方法不同。求维数要处理分子分母的共同消失；求单个权重数则可在负根锥中展开倒数。张量分解最后归结为乘法和最高项的逐次消去。

\subsection{Weyl 维数公式}
\label{b5:subsec:1905:01}

\begin{theorem}[Weyl 维数公式]\label{b5:thm:weyl-dimension}
设 \(\mathfrak g\) 为有限维复半单 Lie 代数，\(\mathfrak h\) 为其 Cartan 子代数，\(\Phi^+\) 为相应根系的一组正根，\(P^+\) 为支配整权集。令 \(\rho=\dfrac12\sum_{\alpha\in\Phi^+}\alpha\)，以 \(h_\alpha\) 表示根 \(\alpha\) 的余根，\((\, ,\,)\) 表示 Killing 型在 \(\mathfrak h^*\) 上诱导的型。对 \(\lambda\in P^+\)，相应有限维简单最高权模 \(L(\lambda)\) 满足
\begin{equation}\label{b5:eq:weyl-dimension}
\dim L(\lambda)
=\prod_{\alpha>0}\frac{(\lambda+\rho,\alpha)}{(\rho,\alpha)}
=\prod_{\alpha>0}\frac{(\lambda+\rho)(h_\alpha)}{\rho(h_\alpha)}.
\end{equation}
\end{theorem}
\begin{proof}
令 \(P\) 为权格，记 \(N=|\Phi^+|\)，并令 \(E=\operatorname{span}_{\mathbb R}\Phi\)、
\(\mathfrak h_{\mathbb R}=\operatorname{span}_{\mathbb R}\{h_\alpha:\alpha\in\Phi\}\)。
根与权在 \(\mathfrak h_{\mathbb R}\) 上取实值。对 \(H\in\mathfrak h_{\mathbb R}\)，作形式代换
\(e^\mu\mapsto\exp\bigl(t\cdot\mu(H)\bigr)\in\mathbb R[\![t]\!]\)，
这给出权格群环 \(\mathbb Z[P]\) 到 \(\mathbb R[\![t]\!]\) 的环同态。对 \(\nu\in P\)，记代换后的交错和为
\(A_\nu(tH)=\sum_w\varepsilon(w)\exp\bigl(t(w\nu)(H)\bigr)\)。其系数为
\[
p_j(H,\nu)=[t^j]A_\nu(tH)
=\frac1{j!}\sum_{w\in W}\varepsilon(w)\bigl((w\nu)(H)\bigr)^j.
\]
右侧把 \(p_j\) 延拓为实空间 \(\mathfrak h_{\mathbb R}\times E\) 上的多项式，分别关于 \(H\) 和 \(\nu\) 为 \(j\) 次齐次。改变求和指标可见，它在两个变量中各自 Weyl 反对称。

关于 \(H\) 的反对称性使 \(p_j\) 在每个根反射的固定超平面 \(\alpha(H)=0\) 上为零，故每个正根线性因子 \(\alpha(H)\) 都整除它。根系约化，所以不同正根不互为非零标量倍数；相应线性多项式是互不相伴的不可约元。由多项式环的唯一分解性，
\(\prod_{\alpha>0}\alpha(H)\) 整除 \(p_j\)，因而 \(j<N\) 时 \(p_j=0\)。当 \(j=N\) 时，比较关于 \(H\) 的次数得到
\[
p_N(H,\nu)=c(\nu)\prod_{\alpha>0}\alpha(H),
\]
其中 \(c(\nu)\) 是关于 \(\nu\) 的 \(N\) 次齐次多项式。因 \(p_N\) 关于 \(\nu\) 反对称，而右侧的 \(H\) 因子与 \(\nu\) 无关，约去该因子得到
\(c(s_\alpha\nu)=-c(\nu)\)。于是 \(c(\nu)\) 在每个固定超平面 \((\nu,\alpha)=0\) 上为零，故被每个线性因子 \((\nu,\alpha)\) 整除。约化性和唯一分解性再次给出整个乘积整除；该乘积为 \(N\) 次齐次，比较次数得到
\(c(\nu)=c\prod_{\alpha>0}(\nu,\alpha)\)，其中 \(c\) 为常数。因此
\[
[t^N]A_\nu(tH)
=c\prod_{\alpha>0}\alpha(H)\prod_{\alpha>0}(\nu,\alpha),
\]
取 \(H\in\mathfrak h_{\mathbb R}\) 使所有 \(\alpha(H)\ne0\)，并令 \(\nu=\rho\)。分母恒等式给出
\[
A_\rho(tH)
=e^{t\cdot\rho(H)}\prod_{\alpha>0}\bigl(1-e^{-t\cdot\alpha(H)}\bigr)
=t^N\prod_{\alpha>0}\alpha(H)+O(t^{N+1}).
\]
因此 \(c=\prod_{\alpha>0}(\rho,\alpha)^{-1}\)。
将 Weyl 特征标公式的乘法形式
\(A_{\lambda+\rho}=A_\rho\,\operatorname{ch}L(\lambda)\)
作同一代换并比较 \(t^N\) 系数，右端的特征标只贡献常数项
\(\sum_\mu m_\lambda(\mu)=\dim L(\lambda)\)。约去非零的
\(\prod_{\alpha>0}\alpha(H)\)，得到第一式。每一项分子分母同时乘
\(2/(\alpha,\alpha)\)，得到余根形式。
\end{proof}

例如对 \(\mathfrak{sl}_3\)，正余根为 \(h_1,h_2,h_1+h_2\)。
\(\lambda=a\omega_1+b\omega_2\)、\(a,b\in\mathbb Z_{\ge0}\) 时，
\[
\dim L(a\omega_1+b\omega_2)
=\frac{(a+1)(b+1)(a+b+2)}2.
\]
所以自然模、对偶自然模、对称平方和伴随模的维数依次为三、三、六、八。取 \(a=b=m\) 时，维数为 \((m+1)^3\)。

\subsection{权重数的计算}
\label{b5:subsec:1905:02}

\begin{definition}\label{b5:def:kostant-partition}
设 \(\Phi\) 为有限维复半单 Lie 代数的根系，\(\Phi^+\) 为一组正根，\(\alpha_1,\ldots,\alpha_\ell\) 为简单根。令 \(E=\operatorname{span}_{\mathbb R}\Phi\)、\(Q=\sum_i\mathbb Z\alpha_i\)、\(Q_+=\sum_i\mathbb Z_{\ge0}\alpha_i\)。对根格 \(Q\) 中的 \(\beta\)，令
\[
\mathcal P(\beta)
=\#\left\{\,(n_\alpha)_{\alpha>0}\in\mathbb Z_{\ge0}^{|\Phi^+|}
\ \middle|\ \sum_{\alpha>0}n_\alpha\alpha=\beta\,\right\}.
\]
若 \(\beta\notin Q_+\)，规定为零；特别地 \(\mathcal P(0)=1\)。
它称为\term[Kostant-fenchai-hanshu]{Kostant 分拆函数}[Kostant partition function]。对不在根格中的向量也规定其值为零。
\end{definition}
\symidx{\(\mathcal P(\beta)\)：Kostant 分拆函数}

由于正根高度为正，固定 \(\beta\) 时所有 \(n_\alpha\) 都有界，故这个数有限。其形式生成函数为
\[
\prod_{\alpha>0}(1-e^{-\alpha})^{-1}
=\sum_{\beta\in Q_+}\mathcal P(\beta)e^{-\beta}.
\]

\begin{theorem}[Kostant 重数公式]\label{b5:thm:kostant-multiplicity}
设 \(\mathfrak g\) 为有限维复半单 Lie 代数，\(\mathfrak h\) 为其 Cartan 子代数，\(\Phi\) 为相应根系，\(\Phi^+\) 为一组正根，\(P\) 为权格，\(P^+\) 为支配整权集，\(W\) 为 Weyl 群。令 \(\rho=\dfrac12\sum_{\alpha\in\Phi^+}\alpha\)、\(\varepsilon(w)=\det(w|_{\operatorname{span}_{\mathbb R}\Phi})\)，以 \(\mathcal P\) 表示 \(\Phi^+\) 的 Kostant 分拆函数。设 \(\lambda\in P^+\)，\(L(\lambda)\) 为相应有限维简单最高权模，\(\mu\in P\)，\(m_\lambda(\mu)=\dim_{\mathbb C}L(\lambda)_\mu\)。则
\begin{equation}\label{b5:eq:kostant-multiplicity}
m_\lambda(\mu)
=\sum_{w\in W}\varepsilon(w)
\mathcal P\bigl(w(\lambda+\rho)-(\mu+\rho)\bigr).
\end{equation}
\end{theorem}
\begin{proof}
在负锥形式环中，将 Weyl 特征标公式的分母倒数展开：
\[
\operatorname{ch}L(\lambda)
=\sum_{w\in W}\varepsilon(w)e^{w(\lambda+\rho)-\rho}
\sum_{\beta\in Q_+}\mathcal P(\beta)e^{-\beta}.
\]
固定 \(e^\mu\) 的系数，所需 \(\beta\) 恰为
\(w(\lambda+\rho)-(\mu+\rho)\)，于是得到公式。Weyl 群有限，且每个系数的形式乘法有限，因而没有交换发散级数的问题。
\end{proof}

在 \(A_2\) 中，正根为 \(\alpha_1,\alpha_2,\alpha_1+\alpha_2\)。所以
\[
\mathcal P(a\alpha_1+b\alpha_2)
=\begin{cases}
\min(a,b)+1,&a,b\in\mathbb Z_{\ge0},\\
0,&\text{其他情形}.
\end{cases}
\]
因为组合中第三条根可选 \(0,\ldots,\min(a,b)\) 次，选定后前两条根的次数即确定。取 \(m\in\mathbb Z_{\ge0}\)，对
\(\lambda=m(\alpha_1+\alpha_2)=m\rho\)、\(\mu=0\)，Kostant 和中只有单位元项非零：其他轨道权
\(w\bigl((m+1)\rho\bigr)-\rho\) 至少有一个简单根系数为负。因此
\[
m_{m\rho}(0)=\mathcal P(m\alpha_1+m\alpha_2)=m+1.
\]
当 \(m=2\) 时，二十七维简单模有三维零权空间。

\ProseParagraphBreak
Freudenthal 递推适合按较高权逐层计算，Kostant 公式则直接把一个重数写成有限个分拆数的交错和。后者的中间项可能很大，并且发生抵消，不意味着它在所有秩中都更省计算。两者均由本章已经证明的等式得到，可用来互相核验。

\ProseParagraphBreak
例如取 \(\mathfrak{sl}_3\) 的最高权 \(\lambda=2\omega_1+\omega_2\)。若支配权 \(\mu=a\omega_1+b\omega_2\) 满足 \(\mu\le\lambda\)，则
\[
\lambda-\mu
=\frac{5-2a-b}{3}\alpha_1+\frac{4-a-2b}{3}\alpha_2.
\]
令 \(a,b\ge0\)，并要求这两个系数为非负整数，得到全部候选
\(\lambda,2\omega_2,\omega_1\)，相对于 \(\lambda\) 的高度依次为零、一、二。
Killing 型的归一化给出
\((\omega_1,\omega_1)=(\omega_2,\omega_2)=1/9\)、
\((\omega_1,\omega_2)=1/18\)。初值 \(m_\lambda(\lambda)=1\)。在 \(\mu=2\omega_2\) 处，递推和中只有 \(\mu+\alpha_1=\lambda\) 有非零重数，其余向上候选都不在 \(\bigcup_{w\in W}w\{\lambda,2\omega_2,\omega_1\}\) 中。于是
\[
\left(\frac{19}{9}-\frac{13}{9}\right)m_\lambda(2\omega_2)
=2\cdot\frac13,
\qquad m_\lambda(2\omega_2)=1.
\]
在 \(\mu=\omega_1\) 处，三个正根各贡献一个非零项：
\(\mu+\alpha_1=3\omega_1-\omega_2=s_2\lambda\)、
\(\mu+\alpha_2=2\omega_2\)、
\(\mu+\alpha_1+\alpha_2=\lambda\)。第一项虽然不支配，却因 Weyl 不变性有重数一；不能在递推中将它删去。三项内积依次为 \(1/2,1/3,1/2\)，故
\[
\left(\frac{19}{9}-\frac79\right)m_\lambda(\omega_1)
=2\left(\frac12+\frac13+\frac12\right),
\qquad m_\lambda(\omega_1)=2.
\]
\(j\ge2\) 的向上权都不在上述 Weyl 轨道并集中，因此没有遗漏的项。三个轨道的大小依次为六、三、三，求和得
\(\dim L(2\omega_1+\omega_2)=6+3+3\cdot2=15\)，与维数公式一致。

\subsection{特征标乘法与张量分解}
\label{b5:subsec:1905:03}

设 \(\lambda,\mu\in P^+\)。先把
\[
F=\operatorname{ch}L(\lambda)\operatorname{ch}L(\mu)
\]
展开为有限权和。由完全可约性，它是若干简单特征标的非负整数和。
选其支集中对权偏序极大的权 \(\nu\)，该权必为某个简单直和项的最高权，因而支配。令 \(a\) 为 \(e^\nu\) 的系数；减去 \(a\,\operatorname{ch}L(\nu)\)，就删掉相应的全部简单直和项。继续这个过程，每次都使表示的总维数下降，有限步后停止。若同时有多个不可比较的极大权，可按任意次序处理。

\begin{example}\label{b5:exm:sl3-natural-dual-tensor}
设 \(V=\mathbb C^3\) 为 \(\mathfrak{sl}_3\) 自然模。由已算出的特征标，
\[
\operatorname{ch}(V\otimes V^*)
=(x_1+x_2+x_3)(x_1^{-1}+x_2^{-1}+x_3^{-1})
=3+\sum_{i\ne j}\frac{x_i}{x_j}.
\]
减去最高权 \(\omega_1+\omega_2\) 的伴随特征标
\(2+\sum_{i\ne j}x_i/x_j\)，剩下一个常数一。因此
\[
V\otimes V^*\cong L(\omega_1+\omega_2)\oplus L(0).
\]
这与 \(\operatorname{End}(V)=\mathfrak{sl}_3\oplus\mathbb CI\) 一致。
类似地，
\[
V\otimes V\cong L(2\omega_1)\oplus L(\omega_2),
\]
因为 \((x_1+x_2+x_3)^2\) 分成
\(\sum_{i\le j}x_ix_j\) 与 \(\sum_{i<j}x_ix_j\)，它们分别是对称平方和外平方的特征标。
\end{example}

\subsection{形式代换与最低次项比较}
\label{b5:subsec:1905:04}

维数公式中，不能直接令所有 \(e^\mu=1\) 后把 Weyl 商的分子与分母相除：两者都会变为零。正确做法是保留一个形式参数 \(t\)，找到共同的首个非零次数，再比较系数。这个次数是正根数 \(N\)，原因是反对称多项式必须含有每一条反射超平面的线性因子。

\ProseParagraphBreak
在 \(A_1\) 中，这个过程只有一次消失。取 \(\omega(H)=1\)，则
\[
A_{(m+1)\omega}(tH)=2(m+1)t+O(t^3),\qquad
A_\omega(tH)=2t+O(t^3),
\]
比值的常数项为 \(m+1\)。在 \(A_2\) 中，三条正根使分子分母都从 \(t^3\) 开始，最低次项的系数之比就是
\((a+1)(b+1)(a+b+2)/2\)。这里的 \(O(t^r)\) 只表示形式幂级数中次数至少为 \(r\) 的项，不调用解析极限。

\ProseParagraphBreak
维数、权重数和张量重数分别回答“总共有多少个向量方向”“某个 Cartan 权出现几次”“某个简单模出现几份”。三个数应分开记录；本章通过特征标把它们联系起来，却没有把它们混作一种重数。
